\documentclass[a4paper, serif, 10pt]{moderncv}

%% moderncv
% style and color
\moderncvstyle{banking}
\moderncvcolor{black}

%% page layout/margins
\usepackage[top=2.5cm, bottom=2.5cm, left=1.5cm, right=1.5cm]{geometry}

\nopagenumbers{}

%% Babel config for German language
% ref:
% - https://latex3.github.io/babel/guides/locale-german.html
\usepackage[ngerman]{babel}

%% fontspec
\usepackage{fontspec}

\IfFontExistsTF{Linux Libertine}{
  \setmainfont[Ligatures={TeX, Common}, Numbers=OldStyle]{Linux Libertine}
}{}
\IfFontExistsTF{Linux Libertine O}{
  \setmainfont[Ligatures={TeX, Common}, Numbers=OldStyle]{Linux Libertine O}
}{}

%% CTeX
% CJK support, used to show CJK characters in the resume
%
% - fontset=none: disable builtin fontset but instead set the CJK font manually
% - heading=false: disable ctex heading
% - punct=kaiming: use kaiming punctuations styles for CJK
% - scheme=plain: use plain scheme, do not override \normalsize font size
% - space=auto: space settings for CJK characters
%
% ref:
% - http://ctan.mirrorcatalogs.com/language/chinese/ctex/ctex.pdf
\usepackage[UTF8, heading=false, punct=kaiming, scheme=plain, space=auto]{ctex}

\IfFontExistsTF{Noto Serif CJK SC}{
  \setCJKmainfont{Noto Serif CJK SC}
}{}
\IfFontExistsTF{Noto Sans CJK SC}{
  \setCJKsansfont{Noto Sans CJK SC}
}{}

%% line spacing
\usepackage{setspace}
\setstretch{1.125}

%% URL styling - use normal text instead of monospace
\urlstyle{same}

% Auto-underline all links
% Defer until \begin{document} so hyperref is already loaded
\AtBeginDocument{%
  \let\oldhref\href
  \renewcommand{\href}[2]{\oldhref{#1}{\underline{#2}}}%
}

%% Patch \httplink and \httpslink to handle full URLs
% The original moderncv macros always prepend http:// or https:// to URLs.
% This causes issues when the URL already has a protocol (e.g., https://example.com)
% resulting in malformed URLs like https://https://example.com.
% This patch checks if the URL already contains "://" and uses it directly if so.
% Note: We use \str_set:Nx (x-expansion) to expand macros like \@homepage first.
\ExplSyntaxOn
\str_new:N \g_yamlresume_url_str

\RenewDocumentCommand{\httpslink}{O{}m}{
  \str_set:Nx \g_yamlresume_url_str {#2}
  \str_if_in:NnTF \g_yamlresume_url_str {://}
    { \url{#2} }
    { \url{https://#2} }
}

\RenewDocumentCommand{\httplink}{O{}m}{
  \str_set:Nx \g_yamlresume_url_str {#2}
  \str_if_in:NnTF \g_yamlresume_url_str {://}
    { \url{#2} }
    { \url{http://#2} }
}
\ExplSyntaxOff

\name{Julia Brandt}{}
\title{Senior Produktmanagerin für digitale B2B-Plattformen}
\phone[mobile]{+49 151 23456789}
\email{julia.brandt@example.de}
\homepage{https://juliabrandt.de}

\address{Torstraße 101, Berlin, Berlin, Deutschland, 10119}{}{}

\extrainfo{{\small \faLinkedin}\ \href{https://www.linkedin.com/in/juliabrandt}{@juliabrandt} {} {} {} • {} {} {}
{\small \faMedium}\ \href{https://medium.com/@juliabrandt}{@@juliabrandt}}

\begin{document}

\maketitle

\section{Basisinformationen}

\cvline{}{\begin{itemize}
\item Produktmanagerin mit über 8 Jahren Erfahrung in der Entwicklung und Skalierung digitaler B2B-Produkte
\item Ausgewiesene Expertise in Produktstrategie, agiler Entwicklung und datengetriebener Entscheidungsfindung
\item Erfolgreiche Führung cross-funktionaler Teams und enge Zusammenarbeit mit Vertrieb, Engineering und Design
\item Leidenschaft für nutzerzentriertes Design und die Umsetzung messbarer Geschäftsziele
\end{itemize}}
\section{Ausbildung}

\cventry{Okt. 2013–März 2016}
        {Master, Management and Technology, Punkte: 1,7}
        {Technische Universität München}
        {\url{https://www.tum.de/}}
        {}
        {\begin{itemize}
\item Vertiefung in Produkt- und Innovationsmanagement sowie quantitative Methoden
\item Masterarbeit über Erfolgsfaktoren digitaler B2B-Plattformen
\end{itemize}
\textbf{Kurse}: Strategisches Management, Innovations- und Technologiemanagement, Empirische Forschungsmethoden, Produktentwicklung}

\cventry{Okt. 2010–Sept. 2013}
        {Bachelor, Wirtschaftsinformatik, Punkte: 2,1}
        {Hochschule für Wirtschaft und Recht Berlin}
        {\url{https://www.hwr-berlin.de/}}
        {}
        {\begin{itemize}
\item Fundierte Kenntnisse in Softwareentwicklung, Datenmodellierung und Geschäftsprozessen
\item Bachelorarbeit zur Einführung agiler Methoden in mittelständischen Unternehmen
\end{itemize}
\textbf{Kurse}: Softwaretechnik, Datenbanken, Business Process Management, Marketing}
\section{Berufserfahrung}

\cventry{Jan. 2022–Heute}
        {Senior Produktmanagerin}
        {FinTech Solutions GmbH}
        {\url{https://www.fintech-solutions.de}}
        {}
        {\begin{itemize}
\item Verantwortlich für Produktstrategie und Roadmap der B2B-Zahlungsplattform mit über 50.000 Nutzern
\item Führung eines cross-funktionalen Teams aus acht Entwicklern, Designern und Data Scientists
\item Steigerung der Aktivierungsrate um 25 \% durch Einführung eines personalisierten Onboardings
\item Durchführung von über 100 Nutzerinterviews und kontinuierlichen A/B-Tests zur Verbesserung der User Experience
\item Enge Abstimmung mit Vertrieb und Customer Success zur Identifikation von Kundenbedürfnissen
\end{itemize}
\textbf{Schlüsselwörter}: Produktstrategie, Roadmapping, Agile Methoden, Datenanalyse, B2B Plattform}

\cventry{März 2019–Dez. 2021}
        {Produktmanagerin}
        {E-Commerce AG}
        {\url{https://www.ecommerce.ag}}
        {}
        {\begin{itemize}
\item Entwicklung und Skalierung der Mobile-Shopping-App von 10.000 auf 250.000 monatliche Nutzer
\item Definition von Produktanforderungen und User Stories in Zusammenarbeit mit internen Stakeholdern
\item Aufbau eines Feedback-Systems zur systematischen Auswertung von Nutzerbewertungen
\item Einführung von SQL-basierten Dashboards zur Überwachung von KPIs und Conversion-Raten
\end{itemize}
\textbf{Schlüsselwörter}: Mobile App, User Stories, KPI-Analyse, Stakeholder-Management}

\cventry{Juni 2016–Feb. 2019}
        {Junior Produktmanagerin}
        {DigitalAgentur GmbH}
        {\url{https://www.digitalagentur.de}}
        {}
        {\begin{itemize}
\item Unterstützung der Produktverantwortlichen bei der Konzeption und Umsetzung von Webprojekten
\item Erstellung von Wireframes und Prototypen für kundenspezifische Lösungen
\item Koordination von agilen Sprints und regelmäßiges Reporting an die Geschäftsleitung
\end{itemize}
\textbf{Schlüsselwörter}: Prototyping, Projektkoordination, Scrum}
\section{Sprachen}

\cvline{Deutsch}{Muttersprache oder zweisprachig \hfill \textbf{Schlüsselwörter}: Muttersprache}
\cvline{Englisch}{Verhandlungssichere Kenntnisse \hfill \textbf{Schlüsselwörter}: Verhandlungssicher, TOEFL 110}
\cvline{Französisch}{Gute Kenntnisse \hfill \textbf{Schlüsselwörter}: Grundkenntnisse}
\section{Fähigkeiten}

\cvline{Produktmanagement}{Experte \hfill \textbf{Schlüsselwörter}: Produktstrategie, Roadmapping, User Stories, Priorisierung, Stakeholder-Management}
\cvline{Agile Methoden}{Experte \hfill \textbf{Schlüsselwörter}: Scrum, Kanban, OKRs, Design Sprint}
\cvline{Datenanalyse}{Erfahren \hfill \textbf{Schlüsselwörter}: SQL, Google Analytics, Amplitude, A/B-Tests}
\cvline{UX/UI Design}{Fortgeschritten \hfill \textbf{Schlüsselwörter}: Figma, Wireframes, Prototyping, Usability Tests}
\section{Auszeichnungen}

\cventry{Jan. 2024}
        {FinTech Solutions GmbH}
        {Award für beste Produktinnovation}
        {}
        {}
        {Auszeichnung für die erfolgreiche Einführung einer KI-basierten Betrugserkennung in der Zahlungsplattform.}
\section{Zertifikate}

\cventry{Juni 2019}
        {Scrum Alliance}
        {Certified Scrum Product Owner (CSPO)}
        {\url{https://www.scrumalliance.org/}}
        {}
        {}

\cventry{Nov. 2020}
        {Product School}
        {Professional Product Discovery}
        {\url{https://productschool.com/}}
        {}
        {}
\section{Veröffentlichungen}

\cventry{Apr. 2023}
        {So gelingt die digitale Transformation im B2B-Vertrieb}
        {Medium}
        {\url{https://medium.com/@juliabrandt}}
        {}
        {\begin{itemize}
\item Praxisbericht über agile Produktentwicklungsmethoden und Kundenorientierung
\item Teilen von Learnings aus der Einführung einer B2B-Plattform
\end{itemize}}
\section{Projekte}

\cventry{Jan. 2023–Jan. 2023}
        {Prototyp eines KI-Modells zur dynamischen Preisgestaltung für E-Commerce-Produkte.}
        {KI-gestützte Preisoptimierung}
        {\url{https://github.com/juliabrandt/preisoptimierung}}
        {}
        {\begin{itemize}
\item Entwicklung eines Machine-Learning-Modells mit Python zur Analyse von Preiselastizitäten
\item Integration in ein Dashboard zur Entscheidungsunterstützung
\end{itemize}
\textbf{Schlüsselwörter}: Machine Learning, Python, Pricing, E-Commerce}
\section{Interessen}

\cvline{Digitale Bildung}{E-Learning, Mentoring, EdTech}
\cvline{Sport}{Laufen, Yoga, Bouldern}
\section{Ehrenamtliche Tätigkeiten}

\cventry{Jan. 2021–Jan. 2023}
        {Mentorin}
        {Frauen in der Digitalwirtschaft e.V.}
        {\url{https://www.fid-ev.de/}}
        {}
        {\begin{itemize}
\item Begleitung von Studentinnen und Berufseinsteigerinnen in die Produktmanagement-Rolle
\item Durchführung von monatlichen Workshops zu Karriereplanung und Verhandlungsstrategien
\end{itemize}}

\end{document}